\chapter{Abenteuer erleben}
\label{chap:adventures}

Neben Kämpfen erleben Helden viele weitere Situationen. Sie schleichen an Wachen vorbei, erklimmen steile Felswände, überwinden Hindernisse oder unterstützen sich gegenseitig bei schwierigen Aufgaben.

Dieses Kapitel beschreibt die Regeln für Zusammenarbeit und Rast. Wie genau eine gemeinsame Handlung in der jeweiligen Situation umgesetzt wird, entscheidet die Spielleitung gemeinsam mit der Gruppe.

\section{Zusammenarbeit}

Helden können Herausforderungen entweder gemeinsam bewältigen oder sich gegenseitig unterstützen.

\subsection{Gruppenproben}

Müssen mehrere Helden gemeinsam dieselbe Herausforderung meistern, kann die Spielleitung eine Gruppenprobe verlangen. Sie bestimmt, welche Helden beteiligt sind und welche Proben zur Situation passen.

\begin{kurzregel}[title=Gruppenprobe]
Die Gesamtschwierigkeit entspricht der regulären Schwierigkeit multipliziert mit der Anzahl der beteiligten Helden.

Alle beteiligten Helden legen die von der Spielleitung bestimmten Proben ab und addieren ihre erzielten \star{}. Erreichen oder überschreiten sie gemeinsam die Gesamtschwierigkeit, ist die Gruppenprobe erfolgreich.
\end{kurzregel}

\begin{beispiel}[title=Beispiel: Gemeinsam schleichen]
Vier Helden möchten sich gemeinsam an einer Wache vorbeischleichen. Die Schwierigkeit beträgt 4.

Da vier Helden beteiligt sind, beträgt die Gesamtschwierigkeit 16\star. Alle vier Helden führen eine passende Schleichen-Probe durch. Erzielen sie gemeinsam mindestens 16\star, bleibt die Gruppe unentdeckt.
\end{beispiel}

\subsection{Helfen}

Ein Held kann einen anderen bei einer Probe unterstützen. Dazu beschreibt der Spieler, wie sein Held helfen möchte. Die Spielleitung entscheidet, ob die Unterstützung sinnvoll ist und welche Probe dafür abgelegt wird.

Mehrere Helden können helfen, wenn jeder von ihnen auf eine eigene sinnvolle Weise zur Lösung beiträgt. Ob eine bestimmte Unterstützung möglich ist, entscheidet die Spielleitung.

\begin{kurzregel}[title=Bei einer Probe helfen]
Der Helfer legt zuerst eine passende Probe ab. Für jeden erzielten \star{} erhält der aktive Held einen zusätzlichen \blueDie{} für seine Probe.

Erzielt der Helfer keinen \star{}, erhält der aktive Held keinen Bonus. Weitere Folgen einer misslungenen Unterstützung bestimmt die Spielleitung anhand der Situation.
\end{kurzregel}

\begin{beispiel}[title=Beispiel: Eine Steintür öffnen]
Gadhra hilft Kiera dabei, eine schwere Steintür zu öffnen. Sie beschreibt, wie sie die Tür mit einem Balken anhebt. Die Spielleitung verlangt dafür eine Attributsprobe mit Kraft.

Gadhra erzielt 2\star. Kiera erhält dadurch 2 zusätzliche \blueDie{} für ihre Probe zum Öffnen der Tür.
\end{beispiel}

\section{Rast}

Helden können rasten, wenn sie ausreichend Zeit und Sicherheit dafür haben. Ob und wann eine Rast möglich ist, entscheidet die Spielleitung.

Jeder Held rastet unabhängig von den anderen. Während einige Helden rasten, können andere beispielsweise Wache halten oder anderen Tätigkeiten nachgehen.

Während einer Rast sind nur leichte Tätigkeiten möglich. Kämpfen, Reisen, anstrengende Proben oder vergleichbare Belastungen unterbrechen die Rast. Im Zweifel entscheidet die Spielleitung, ob eine Tätigkeit noch als Rast gilt.

\begin{kurzregel}[title=Eine Stunde rasten]
Für jede vollständig abgeschlossene Stunde Rast erhält ein Held:

\begin{itemize}
  \item Lebenspunkte in Höhe seines KON-Wertes, höchstens bis zu seinem Maximum,
  \item Mana in Höhe seines CHA-Wertes, höchstens bis zu seinem Maximum,
  \item alle erschöpften Karten wieder bereitgemacht.
\end{itemize}
\end{kurzregel}

Wird eine Rast vor Ablauf einer vollständigen Stunde unterbrochen, erhält der Held für diese angefangene Stunde keine Vorteile.

Nach einer vollständig abgeschlossenen Stunde werden alle umgedrehten Karten -- beispielsweise Talente, Zauber und Ausrüstung -- wieder aufgedeckt, sofern eine Regel nichts anderes bestimmt.
